\section{Conclusion}
\label{sec:conclusion}

We began by observing that on 18\% of factuality questions, token entropy and answer consistency point in opposite directions---and each method is right about its own subset. This counterintuitive finding motivated us to systematically investigate whether these uncertainty signals capture fundamentally different hallucination patterns.

In this work, we conducted the first systematic comparison of entropy-based and consistency-based uncertainty quantification for LLM hallucination detection on the same benchmark. Our key insight is that token entropy captures epistemic uncertainty while N-sample consistency captures generation stability---orthogonal phenomena measuring different aspects of model uncertainty.

Our main contributions are:

\begin{enumerate}
    \item \textbf{Empirical demonstration of orthogonality:} Entropy and consistency share only 5\% of variance ($r = 0.228$).

    \item \textbf{Mechanism validation:} Consistency shows a large effect size (Cohen's $d = 1.068$) distinguishing correct from incorrect responses.

    \item \textbf{Complementarity analysis:} 18.1\% of questions exhibit method disagreement, with each method achieving AUROC $> 0.75$ on its winning subset.
\end{enumerate}

As language models become more capable yet continue to hallucinate, understanding the orthogonal structure of uncertainty signals becomes critical. Our finding that entropy and consistency capture complementary failure modes suggests that robust hallucination detection requires multiple lenses---no single signal suffices. We hope this work encourages the community to move beyond single-signal approaches toward principled multi-signal frameworks for trustworthy AI systems.
